% frontmatter/abstract.tex
\chapter*{Abstract}
\addcontentsline{toc}{chapter}{Abstract}
This report documents and evaluates, through a documentary reverse-engineering approach over the
repository, the design and implementation of \textbf{Democra}, a multi-tenant Software as a Service
(SaaS) web platform for the comprehensive management of volunteer organizations (NGOs). The system
comprises a Single Page Application (SPA) client built with React, TypeScript and Vite, and a REST
API implemented with Node.js and Express on top of the Supabase data platform (PostgreSQL 16). At the
data level, it adopts a multi-schema model with tenant isolation enforced through Row Level Security
(RLS) and database functions. Security relies on Supabase authentication, Role-Based Access Control
(RBAC) with a deny-by-default policy, a risk evaluation engine with two-step verification (OTP),
forensic auditing, and soft deletion with retention policies. Methodologically, the development is governed
by the DDS methodology (Document-Driven SDLC assisted by AI), organized into eight lifecycle
phases. This methodology operationalizes the document-as-gate practice established across the
technology industry through design docs and RFCs (e.g., Amazon's Working Backwards, Google, Uber),
grounded in Parnas and Clements' rational design process. The research is applied-technological and descriptive, with a non-experimental cross-sectional
design. The study variables are Functioning and Usability; their indicators are validated through
documentary code analysis, automated testing (Jest, Vitest, supertest, pgTAP) and code coverage and automated testing. Quantitative coverage and performance results are structured and
reserved for execution in a controlled environment, without reporting unverified values.

\vspace{0.3cm}
\noindent\textbf{Keywords:} multi-tenant SaaS, NGO, volunteering, React, Supabase, PostgreSQL, RBAC,
RLS, quality assurance, DDS methodology.
